\section{What the model reveals in practice}
\label{sec:examples}

The examples answer three practical questions: how local skew arises from the
law, what an endpoint coefficient says about unquoted risk, and why a close
vanilla fit does not identify a fine-grained density.  They use frozen SPY and
NVDA option chains from the 2026-08-03 US session, with eight expiries per
underlier.  The fits use the exponential-tail subclass
$\alpha_-=\alpha_+=0$, order $N=16$ subject to
\eqref{eq:orderguard}, haircut bands, and ridge
$\omega_{\rm reg}=10^{-6}$ with $\nu=1$.  The examples therefore test the
original LQD subclass; the generalized tail exponent is evaluated by separate
scenario fits.

\subsection{A liquid index: skew is a digital imbalance}

\begin{figure}[htbp]
\centering
\paperfig{0.94\textwidth}{fig_spy_node.pdf}
\caption{The SPY December 2026 slice.  Left: quoted volatilities and the
fitted smile.  Center: signed residuals in volatility basis points.  Right:
the fitted density and an ATM-matched normal density.  The left-shifted mass
and richer model digital produce the negative equity skew.}
\label{fig:spy_node}
\end{figure}

The slice contains \MacSpyDecNQuotes\ quotes and fits them to
\MacSpyDecRmsBp\ volatility basis points rms.  The exact ATM formulas give
volatility \MacSpyDecAtmVolPct\%, skew $\MacSpyDecSkew$, and curvature
$\MacSpyDecCurv$.  The forward rank is \MacSpyDecForwardRankPct\%, below the
matched Black rank, so $\Delta<0$ and \eqref{eq:atm_sigma1} gives a negative
skew.  The log-contract variance corresponds to
\MacSpyDecVarSwapPct\% annualized volatility.

For the listed $\MacTicketStrike$ call, $k=\MacTicketK$ and the strike rank is
\MacTicketRankPct\%.  The asset-share and cash entries are
$\MacTicketShare$ and $\MacTicketCash$, hence
\[
c(\MacTicketK)=\MacTicketShare-\MacTicketCash=\MacTicketPrice.
\]
This is \MacTicketDollarPrice\ per share and corresponds to
\MacTicketIvPct\% implied volatility.

The endpoint coefficients of the exponential subclass are
$\lambda_-=\MacSpyLamL$ and $\lambda_+=\MacSpyLamR$.  They imply Lee slopes
$\beta_L=\MacSpyBetaL$ and $\beta_R=\MacSpyBetaR$.  These are asymptotic
properties of the fitted continuation, not measurements at the outer quoted
strikes.

\subsection{A single name: vanilla fit does not settle tail policy}

The long NVDA slice has endpoint coefficients
$\lambda_-=\MacNvdaLamL$ and $\lambda_+=\MacNvdaLamR$.  Its left coefficient
exceeds one.  This is admissible because negative log-returns do not threaten
the forward; it implies that $\E[Y^{-1}]$ is infinite in the exponential-tail
subclass.  A scenario with $\alpha_->0$ keeps the fitted central law but makes
all inverse moments finite.  Prices of inverse-sensitive or crash-tail
positions distinguish the two policies more directly than vanilla fit error.

On the shortest NVDA expiry, \MacGuardQuoteCountZeroDte\ retained quotes give
$N_{\rm eff}=\MacGuardEffOrderZeroDte$.  The order guard prevents a sparse
strip from assigning an independent high-frequency shape mode to each quote.

\subsection{Basis order: the smile may agree while the density differs}

The synthetic benchmark is a martingale-normalized mixture of two lognormal
components.  Its density is bimodal, while its implied-volatility curve has
only a mild shoulder.

\begin{figure}[htbp]
\centering
\paperfig{0.9\textwidth}{fig_order_control.pdf}
\caption{The same bimodal target fitted at orders 6 and 16.  Left: the two
implied-volatility curves nearly coincide over the quoted range.  Right: the
densities differ; the lower-order model shifts the modes and fills part of the
trough.  Vanilla residuals alone do not identify fine density structure.}
\label{fig:order_control}
\end{figure}

At order 16 the $L^1$ density error is \MacHumpLOneNSixteen; at order 6 it is
\MacHumpLOneNSix.  Yet the implied-volatility curves differ by at most
\MacHumpIvAgreeBp\ volatility basis points on the comparison grid.  Basis
order should therefore be selected for the downstream use.  Broad vanilla
prices tolerate a lower order; digitals and narrow butterflies require a
finer rank representation and stronger identification checks.

\subsection{Relation to other smile models}

Breeden--Litzenberger identify the call curve with a density
\citep{BreedenLitzenberger1978}.  Quantile modeling supplies the rank
coordinate \citep{Parzen1979,Gilchrist2000}; Petersen--M\"uller use a
log-quantile-density transform for density-valued data
\citep{PetersenMuller2016}; the metalog family is a related direct quantile
expansion \citep{Keelin2016}.  Lee and Benaim--Friz connect moments to
implied-variance wings \citep{Lee2004,BenaimFriz2009}.

SVI and SSVI specify total variance directly and characterize admissible
parameter regions \citep{Gatheral2004,GatheralJacquier2014}.  Constrained
price smoothers work directly with convex call curves \citep{Fengler2009}.
LQD instead makes the law primary.  Its main computational object is the
pair $(x,G)$: the transport supplies density and the ledger supplies prices
and calendar order.

\subsection{What follows, what is fitted, and what is chosen}
\label{sec:limits}

\begin{table}[htbp]
\centering\small
\begin{tabular}{@{}>{\raggedright\arraybackslash}p{0.30\textwidth}
>{\raggedright\arraybackslash}p{0.31\textwidth}
>{\raggedright\arraybackslash}p{0.31\textwidth}@{}}
\toprule
\textbf{Derived} & \textbf{Fitted and checked} & \textbf{Chosen}\\
\midrule
Positive density, unit mass, unit forward, decreasing convex calls, and
butterfly-free prices & Quote residuals, interpolation monotonicity,
Durrleman margin, and grid convergence & Basis order, ridge, vega floor,
quote-band rule, and numerical tolerances\\[4pt]
Tail rate \eqref{eq:righttailrate}, moment domain \eqref{eq:strip}, Lee slopes
or sublinear wing law & Tail-quadrature accuracy, operational tail guard, and
scenario sensitivity beyond quoted ranks & Tail exponents $\alpha_\pm$ and
any additional endpoint-moment limits\\[4pt]
Prices, digitals, density, ATM handles, log-contract variance, and analytic
Jacobian from $(x,G)$ & Multi-start stability, active calendar constraints,
and the full-line quantile-root certificate & Common or monotone tail policy
across expiry and the permitted quote-band relaxation\\[4pt]
Full-line call order, ledger order, and noncrossing total variance are
equivalent & The accepted maturity stack satisfies the order to its stated
normalized-price tolerance & Dynamics between the fitted marginal laws are
not selected by the surface\\
\bottomrule
\end{tabular}
\caption{The LQD operating contract.  The first column follows from the
model; the second is measured on the implementation; the third is supplied
before calibration.}
\label{tab:contract}
\end{table}

A fitted surface is published with its quote residuals, effective orders,
quoted-rank spans, tail exponents and coefficients, moment limits, numerical
margins, and minimum adjacent ledger gap.  These quantities are sufficient to
reconstruct how the surface is fitted, extrapolated, and certified.
